% ---- pg-401 （书页 389） ----

对于极大似然法, 第一因子表现为一般运动能力因子, 但是载荷的模式不如主成分因子分析解那么有说服力. 第二因子本质上是强壮因子, 因为铅球和铁饼两个项目在此因子上有大的载荷. 第三因子是跑的持续力因子, 因为 400 米跑和 1 500 米跑两个项目在此因子上有大的载荷. 第四因子也是不易解释, 似乎与跳跃项目有关, 或可称为腿部力量. 我们将在讨论过因子旋转后, 在例 9.11 中回到本节的因子解释问题.

四因子的主成分解占有(标准化的)总样本方差的绝大部分, 尽管某些项目(如标枪)的特殊方差比较大. 这意味着某些项目需要其他项目所不需要的唯一的或特殊的能力. 四因子的极大似然解占有的总样本方差没有那么多, 但是, 如同下面的残差矩阵所显示, 极大似然估计的 $\hat{\mathbf{L}}$ 和 $\hat{\boldsymbol{\Psi}}$ 在再生 $\mathbf{R}$ 时比主成分估计的 $\mathbf{L}$ 和 $\boldsymbol{\Psi}$ 效果更好.

主成分:
\[
\mvfit{$\displaystyle
{\small
\setlength{\arraycolsep}{5pt}
\mathbf{R}-\tilde{\mathbf{L}}\tilde{\mathbf{L}}'-\tilde{\boldsymbol{\Psi}}=
\left[\begin{array}{rrrrrrrrrr}
0&-0.082&-0.006&-0.021&-0.068&0.031&-0.016&0.003&0.039&0.062\\
-0.082&0&-0.046&0.033&-0.107&-0.078&-0.048&-0.059&0.042&0.006\\
-0.006&-0.046&0&0.006&-0.010&-0.014&-0.003&-0.013&-0.151&0.055\\
-0.021&0.033&0.006&0&-0.038&-0.204&-0.015&-0.078&-0.064&-0.086\\
-0.068&-0.107&-0.010&-0.038&0&0.096&0.025&-0.006&0.030&-0.074\\
0.031&-0.078&-0.014&-0.204&0.096&0&0.015&-0.124&0.119&0.085\\
-0.016&-0.048&-0.003&-0.015&0.025&0.015&0&-0.029&-0.210&0.064\\
0.003&-0.059&-0.013&-0.078&-0.006&-0.124&-0.029&0&-0.026&-0.084\\
0.039&0.042&-0.151&-0.064&0.030&0.119&-0.210&-0.026&0&-0.078\\
0.062&0.006&0.055&-0.086&-0.074&0.085&0.064&-0.084&-0.078&0
\end{array}\right]
}
$}
\]

极大似然:
\[
\mvfit{$\displaystyle
{\small
\setlength{\arraycolsep}{5pt}
\mathbf{R}-\hat{\mathbf{L}}\hat{\mathbf{L}}'-\hat{\boldsymbol{\Psi}}=
\left[\begin{array}{rrrrrrrrrr}
0&0.000&0.000&-0.000&-0.000&0.000&-0.000&0.000&-0.001&0.000\\
0.000&0&-0.002&0.023&0.005&0.017&-0.003&-0.030&0.047&-0.024\\
0.000&-0.002&0&0.004&-0.000&-0.009&0.000&-0.001&-0.001&0.000\\
-0.000&0.023&0.004&0&-0.002&-0.030&-0.004&-0.006&-0.042&0.010\\
-0.000&0.005&-0.001&-0.002&0&-0.002&0.001&0.001&0.000&-0.001\\
0.000&-0.017&-0.009&-0.030&-0.002&0&0.022&0.069&0.029&-0.019\\
-0.000&-0.003&0.000&-0.004&0.001&0.022&0&-0.000&-0.000&0.000\\
0.000&-0.030&-0.001&-0.006&0.001&0.069&-0.000&0&0.021&0.011\\
-0.001&0.047&-0.001&-0.042&0.001&0.029&-0.000&0.021&0&-0.003\\
0.000&-0.024&0.000&0.010&-0.001&-0.019&0.000&0.011&-0.003&0
\end{array}\right]
}
$}
\]
\hfill$\blacksquare$

\subsection*{公共因子个数的大样本检验}

正态总体的假设直接引出模型充分性的检验. 设 $m$ 个公共因子的模型成立, 此处 $\boldsymbol{\Sigma}=\mathbf{L}\mathbf{L}'+\boldsymbol{\Psi}$, 并且检验 $m$ 个公共因子模型的充分性等价于检验
\phantomsection\label{eq:9-33}
\begin{equation}
\mvfit{$\displaystyle
H_0:\ \underset{(p\times p)}{\boldsymbol{\Sigma}}=\underset{(p\times m)}{\mathbf{L}}\ \underset{(m\times p)}{\mathbf{L}'}+\underset{(p\times p)}{\boldsymbol{\Psi}}
$} \tag{9-33}
\end{equation}
备选 $H_1$: $\boldsymbol{\Sigma}$ 是任意其他的正定矩阵. 当 $\boldsymbol{\Sigma}$ 没有特殊形式时, 似然函数的最大值[见\hyperref[eq:4-18]{式(4-18)}和
